\chapter{Teoría local de Cauchy y aplicaciones}
\section{Preliminares}
\begin{definition}[Curva]
	Una \textbf{curva} es una aplicación continua $\displaystyle \gamma : \left[a,b\right] \to \R^{n} $. 
\end{definition}
\begin{notation}
	Para evitar problemas asumiremos que $\displaystyle \gamma  $ es $\displaystyle \mathcal{C}^{1} $ y que $\displaystyle \gamma'\neq 0 $ en $\displaystyle \left[a,b\right]  $. 
\end{notation}
Vamos a establecer la siguiente relación de equivalencia entre curvas:
\[\gamma_{1} \sim \gamma_{2} \iff \exists\varphi:\left[c,d\right] \to \left[a,b\right] , \; \varphi' > 0, \; \gamma_{2}=\gamma_{1}\circ \varphi .\]
Podemos observar que $\displaystyle \varphi $ es un difeomorfismo que preserva la orientación de la curva. Así, podemos considerar el conjunto cociente correspondiente obteniendo el \textbf{conjunto de curvas orientadas} $\displaystyle \mathcal{C}^{1} $. Esto nos permite identificar cada curva con su conjunto imagen $\displaystyle \Gamma = \left[\gamma \right]  $. 
\begin{notation}
	Si la orientación es la opuesta escribimos $\displaystyle \Gamma^{-}=\left[\gamma^{-}\right]  $. En particular, dada $\displaystyle \gamma : \left[a,b\right] \to \R^{n} $ tendremos que $\displaystyle \gamma^{-}\left(t\right)=\gamma\left(b+a-t\right) $. 
\end{notation}
\begin{definition}[Curva $\displaystyle \mathcal{C}^1 $ a trozos]
	Decimos que una curva $\displaystyle \gamma : \left[a,b\right] \to \R^{n} $ es $\displaystyle \mathcal{C}^{1} $ \textbf{a trozos} si existe una partición de $\displaystyle \left[a,b\right]  $, $\displaystyle a_{0}= a < a_{1} < \cdots < a_{m} = b $ tal que la restricción de $\displaystyle \gamma  $ a cada subintervalo $\displaystyle \left[a_{j}, a_{j+1}\right]  $ para $\displaystyle j = 0, \ldots, m - 1 $ define una curva parametrizada de clase $\displaystyle \mathcal{C}^{1} $.
\end{definition}

Como vimos en cálculo integral, todo lo anterior se puede generalizar a curvas que sean $\displaystyle \mathcal{C}^{1} $ a trozos. 
\begin{definition}[Longitud de una curva]
	Dada una curva $\displaystyle \gamma : \left[a,b\right] \to \R^{n} $ de clase $\displaystyle \mathcal{C}^{1} $ a trozos, se define la \textbf{longitud} de $\displaystyle \gamma  $ como
	\[\ell\left(\gamma \right)=\int^{b}_{a} \|\gamma'\left(t\right)\| \; dt .\]
\end{definition}
Es sencillo comprobar que si $\displaystyle \gamma  $ y $\displaystyle \sigma  $ son curvas equivalentes, entonces $\displaystyle \ell\left(\gamma \right)=\ell\left(\sigma \right) $, también $\displaystyle \ell\left(\gamma \right)=\ell\left(\gamma^{-} \right) $. De esta forma, para cada curva orientada $\displaystyle \Gamma  $ de clase $\displaystyle \mathcal{C}^{1} $ a trozos puede definirse $\displaystyle \ell\left(\Gamma \right)=\ell\left(\gamma \right) $ con $\displaystyle \Gamma = \left[\gamma \right]  $. 
Además, $\displaystyle \ell\left(\Gamma \right)=\ell\left(\Gamma^{-}\right) $. \\

Una operación que usamos con frecuencia es la \textbf{concatenación de curvas}. Si $\displaystyle \Gamma_{1}, \ldots, \Gamma_{m} $ son curvas de clase $\displaystyle \mathcal{C}^{1} $ a trozos en $\displaystyle \R^{n} $ tales que el punto final de $\displaystyle \Gamma_{j} $ coincide con el punto inicial de $\displaystyle \Gamma_{j + 1} $, se puede definir una curva $\displaystyle \Gamma $ que recorre de manera consecutiva $\displaystyle \Gamma_{1} , \ldots, \Gamma_{m} $ que denotaremos por $\displaystyle \Gamma = \Gamma_{1} \cup \cdots \cup \Gamma_{m} $. 
Para encontrar una parametrización de $\displaystyle \Gamma  $ basta con tomar parametrizaciones $\displaystyle \gamma_{1}, \ldots, \gamma_{m} $ de $\displaystyle \Gamma_{1}, \ldots, \Gamma_{m} $ definidas en intervalos $\displaystyle \left[a_{j}, b_{j}\right]  $ para $\displaystyle j \in \left\{ 1, \ldots, m\right\}  $ de manera que $\displaystyle b_{j}=a_{j+1} $ y que $\displaystyle \gamma_{j}\left(b_{j}\right)=\gamma_{j+1}\left(a_{j+1}\right) $. Así, definimos
\[\gamma : [a_{1}, b_{m}] \to \R^{n} : t \to \gamma_{j}\left(t\right), \; t \in \left[a_{j}, b_{j}\right]  .\]
Claramente $\displaystyle \Gamma $ será la clase de las curvas equivalentes a esta $\displaystyle \gamma  $. 
\begin{definition}
Sea $\displaystyle \Gamma $ una curva orientada de clase $\displaystyle \mathcal{C}^{1} $ a trozos en $\displaystyle \R^{n} $. 
\begin{itemize}
	\item Diremos que $\displaystyle \Gamma $ es una \textbf{curva cerrada} si admite una parametrización $\displaystyle \gamma : \left[a,b\right] \to \R^{n} $ tal que $\displaystyle \gamma\left(a\right)=\gamma\left(b\right) $.
	\item Diremos que $\displaystyle \Gamma $ es una \textbf{curva simple} si admite una parametrización $\displaystyle \gamma : \left[a,b\right] \to \R^{n} $ que es inyectiva.
	\item Diremos que $\displaystyle \Gamma $ es una \textbf{curva cerrada simple} si admite una parametrización $\displaystyle \gamma : \left[a,b\right] \to \R^{n} $ tal que $\displaystyle \gamma\left(a\right)=\gamma\left(b\right) $ y la restricción de $\displaystyle \gamma  $ a $\displaystyle [a,b) $ es inyectiva.
\end{itemize}
\end{definition}
\begin{observation}
Para las curvas paramétricas cerradas $\displaystyle \mathcal{C}^{1} $ a trozos se pueden considerar reparametrizaciones $\displaystyle \gamma : \mathbb{S}^{1} \to \R^{n} $ para evitar problemas. Usando este punto de vista
\[\gamma \sim \sigma \iff \exists \varphi : \mathbb{S}^{1} \to \mathbb{S}^{1} ,\]
homeomorfismo $\displaystyle \mathcal{C}^{1} $ a trozos que preserva orientación tal que $\displaystyle \sigma =\gamma\circ\varphi $. 
\end{observation}
\section{Integración de funciones con valores en $\displaystyle \R^n $}
Sea $\displaystyle \varphi : A \subset \R^{d}\to \R^{n} $, podemos escribir $\displaystyle \varphi = \left(\varphi_{1}, \ldots, \varphi_{n}\right) $ y definir
\[\int _{A} \varphi \; d x : = \left(\int _{A}\varphi_{1} \; dx, \ldots, \int _{A}\varphi_{n} \; dx\right),\]
siempre que las integrales $\displaystyle \int _{A}\varphi_{i} \; dx $ estén bien definidas $\displaystyle \forall i \in \left\{ 1, \ldots, n\right\}  $. 
\begin{prop}[Propiedades de la integral]
Sean $\displaystyle \varphi, \psi : A\subset \R^{d} \to \R^{n} $ integrables y $\displaystyle \alpha, \beta \in \R $. Entonces, 
\begin{enumerate}
\item $\displaystyle \int _{A}\left(\alpha\varphi + \beta\varphi\right) \; dx = \alpha\int _{A}\varphi \; dx + \beta\int _{A}\psi \; dx $. 
\item $\displaystyle \left\|\int _{A}\varphi \; dx\right\| \leq \int _{A}\|\varphi\| \; dx$. 
\item \textbf{Teorema fundamental del cálculo}. Para toda $\displaystyle g : \left[a,b\right] \to \R^{n} $ de clase $\displaystyle \mathcal{C}^{1}\left(\left[a,b\right] \right) $ se tiene que 
	\[g\left(t\right)=g\left(a\right)+\int^{t}_{a} g'\left(x\right) \; dx .\]
\item Si $\displaystyle A = A_{1} \cup A_{2} $ con $\displaystyle m\left(A_{1} \cap A_{2}\right)=0 $, entonces $\displaystyle \int _{A}\varphi \; dx = \int _{A_{1}}\varphi \; dx + \int _{A_{2}}\varphi \; dx $. 
\end{enumerate}
\end{prop}
\begin{proof}
Las propiedades \textbf{(1)}, \textbf{(3)} y \textbf{(4)} se deducen fácilmente a partir de sus resultados correspondientes para funciones con valores en $\displaystyle \R $. Para demostrar \textbf{(2)} denotamos 
\[x = \left(x_{1}, \ldots, x_{n}\right)=\left(\int _{A}\varphi_{1} , \ldots, \int _{A}\varphi_{n} \right)=\int _{A}\varphi .\]
Podemos suponer que $\displaystyle x \neq 0 $. Usando la linealidad de la integral de funciones con valores en $\displaystyle \R $ y la desigualdad Cauchy-Schwarz en $\displaystyle \R^{n} $,
\[\|x\|^{2}=\left\langle x, x \right\rangle  =\sum^{n}_{j= 1}x_{j}\int _{A}\varphi_{j}\left(t\right) \; dt = \int _{A}\sum^{n}_{j = 1}x_{j}\varphi_{j}\left(t\right) \; dt \leq \int _{A}\|x\| \|\varphi\| \; dt = \|x\|\int _{A}\|\varphi\|,\]
deducimos \textbf{(2)} al dividir por $\displaystyle \|x\| $. 
\end{proof}

\section{Integración de funciones complejas en $\displaystyle \C $}
Al identificar $\displaystyle \R^{2} $ con $\displaystyle \C $ podemos ver que dada $\displaystyle \varphi : \left[a,b\right] \to \C $ con $\displaystyle \varphi\left(t\right)=x\left(t\right) + iy\left(t\right) $ se define de manera natural la \textbf{integral compleja} como
	\[\int^{b}_{a} \varphi\left(t\right) \; dt := \int^{b}_{a} x\left(t\right) \; dt + i\int^{b}_{a} y\left(t\right) \; dt .\]
	Podemos observar que por la propiedad \textbf{(2)} anterior tendremos que 
\[ \left|\int^{b}_{a} \varphi\left(t\right) \; dt\right| \leq \int^{b}_{a} \left|\varphi\left(t\right)\right| \; dt .\]
\begin{definition}[Integral compleja sobre una curva]
	Sea $\displaystyle f : \Omega\subset \C \to \C $ continua y $\displaystyle \gamma : \left[a,b\right] \to \Omega  $ una curva parametrizada de clase $\displaystyle \mathcal{C}^{1} $ a trozos. Se define la \textbf{integral de $\displaystyle f $ sobre $\displaystyle \gamma  $} como
	\[\int _{\gamma }f : = \int^{b}_{a} f\left(\gamma\left(t\right)\right)\gamma'\left(t\right) \; dt .\]
\end{definition}
Si escribimos $\displaystyle f = u + iv $ y $\displaystyle \gamma\left(t\right)=x\left(t\right)+iy\left(t\right) $, obtenemos una expresión equivalente para la integral,
\[
\begin{split}
	\int _{\gamma }f = & \int^{b}_{a} \left(u\left(x\left(t\right), y\left(t\right)\right)x'\left(t\right)-v\left(x\left(t\right), y\left(t\right)\right)y'\left(t\right)\right) \; dt 
\\
			  + & i \int^{b}_{a} \left(u\left(x\left(t\right), y\left(t\right)\right)y'\left(t\right) + v\left(x\left(t\right), y\left(t\right)\right)x'\left(t\right)\right) \; dt..
\end{split}
\]

\begin{observation}
El producto que aparece en la integral anterior \textbf{no es el producto escalar}, es el producto usual definido en $\displaystyle \C $. Esto tiene varias consecuencias:
\begin{enumerate}
\item El valor de $\displaystyle \int _{\gamma }f $ es en general un número complejo por lo que no coincide con el valor de la integral del campo vectorial $\displaystyle f $ a lo largo de $\displaystyle \gamma  $ en $\displaystyle \R^{2} $ (sería un valor real).
\item \textbf{No es cierto en general} que $\displaystyle \left|\int _{\gamma }f \right|\leq \int _{\gamma } \left|f\right| $. Esto no tendría sentido puesto que, como hemos dicho antes, el término de la derecha es en general un número complejo. Ni siquiera tiene por qué cumplirse aún suponiendo que ambos términos son reales, como se verá en un ejemplo a continuación. 
\end{enumerate}
\end{observation}
\begin{eg}
	Cogemos $\displaystyle f\left(z\right)=\frac{1}{z} $ en $\displaystyle \mathbb{S}^{1} $. Observamos que si cogemos la parametrización $\displaystyle \gamma\left(t\right)=e^{it} $ con $\displaystyle t \in \left[0, 2\pi \right]  $ tendremos que 
	\[\int _{\gamma }\frac{1}{z} \; dz = \int^{2\pi }_{0} \frac{1}{e^{it}}ie^{it} \; dt = 2\pi i \Rightarrow  \left|\int _{\gamma }\frac{1}{z} \right|=2\pi.\]
Sin embargo, 
\[\int_{\gamma} \left|f\right| = \int _{\gamma }1 =\int^{2\pi }_{0} 1 \cdot ie^{it} =0 .\]
La desigualdad $\displaystyle \left|\int _{\gamma }f \right|\leq \int _{\gamma } \left|f\right| $ no se cumple puesto que $\displaystyle 2\pi > 0 $. 
\end{eg}

\begin{prop}[Propiedades de la integral compleja]
Sean $\displaystyle f, g : \Omega \subset \C \to \C $ continuas, $\displaystyle \alpha, \beta \in \C $ y $\displaystyle \gamma : \left[a,b\right] \to \Omega  $ curva parametrizada $\displaystyle \mathcal{C}^{1} $ a trozos. Entonces, 
	\begin{enumerate}
	\item $\displaystyle \int _{\gamma }\alpha f + \beta g = \alpha \int_{\gamma } f +\beta \int_{\gamma} g  $.
	\item $\displaystyle \int _{\gamma^{-} }f =-\int _{\gamma }f $. 
	\item $\displaystyle \left|\int _{\gamma }f \right|\leq \sup_{z \in \gamma\left(\left[a,b\right] \right)} \left|f\left(z\right)\right| \ell\left(\gamma \right) $.
	\end{enumerate}
\end{prop}
\begin{proof}
Las propiedades \textbf{(1)} y \textbf{(2)} se vieron en cálculo integral. Demostramos \textbf{(3)}, 
\[
\begin{split}
	\left|\int _{\gamma }f \right|= & \left|\int^{b}_{a} f\left(\gamma\left(t\right)\right)\gamma'\left(t\right) \; dt\right|\leq \int^{b}_{a} \left|f\left(\gamma\left(t\right)\right)\right| \left|\gamma'\left(t\right)\right| \; dt \\
	\leq & \sup_{z \in \gamma\left(\left[a,b\right] \right)} \left|f\left(z\right)\right| \int^{b}_{a} \left|\gamma'\left(t\right)\right| \; dt = \sup_{z \in \gamma\left(\left[a,b\right] \right)} \left|f\left(z\right)\right| \ell\left(\gamma \right).
\end{split}
\]
\end{proof}
\begin{prop}
	Sean $\displaystyle \gamma : \left[a,b\right] \to \Omega  $ y $\displaystyle \sigma : \left[c,d\right] \to \Omega  $ curvas parametrizadas de clase $\displaystyle \mathcal{C}^{1} $ a trozos y supongamos que son equivalentes. Entonces, para toda función continua $\displaystyle f : \Omega \subset \C \to \C $ se tiene que 
	\[\int _{\gamma }f =\int _{\sigma  }f  .\]
\end{prop}
\begin{proof}
	Haremos la demostración suponiendo que $\displaystyle \gamma  $ y $\displaystyle \sigma  $ son $\displaystyle \mathcal{C}^{1} $ puesto que la extensión al caso general es evidente. Supongamos pues que $\displaystyle \phi : \left[c,d\right] \to \left[a,b\right]  $ es biyectiva de clase $\displaystyle \mathcal{C}^{1} $ con $\displaystyle \phi'\left(t\right)>0 $ con $\displaystyle \sigma\left(s\right)=\gamma\left(\phi\left(s\right)\right) $. Aplicando el teorema de cambio de variable tenemos que 
	\[\int _{\gamma }f =\int^{b}_{a} f\left(\gamma\left(t\right)\right)\gamma'\left(t\right) \; dt = \int^{d}_{c} f\left(\gamma\left(\phi\left(s\right)\right)\right)\gamma'\left(\phi\left(s\right)\right)\phi'\left(s\right) \; d s = \int^{d}_{c} f\left(\sigma\left(s\right)\right) \sigma'\left(s\right) \; d s = \int _{\sigma }f .\]
\end{proof}

Debido a este hecho, para cada $\displaystyle \Gamma $ curva orientada de clase $\displaystyle \mathcal{C}^{1} $ a trozos contenida en $\displaystyle \C $ tendremos que podemos definir
	\[\int _{\Gamma}f := \int _{\gamma }f ,\]
con $\displaystyle \gamma  $ parametrización de $\displaystyle \Gamma  $, para cualquier función continua $\displaystyle f : \Gamma \to \C $. Las propiedades de la proposición anterior tienen análogos obvios para $\displaystyle \int _{\Gamma }f $. 
Asímismo, dada una concatenación de curvas $\displaystyle \Gamma = \bigcup^{m}_{j = 1}\Gamma_{j} $, es sencillo ver que
\[\int _{\Gamma}f = \sum^{m}_{j =1}\int _{\Gamma_{j}}f .\]
En el caso particular de $\displaystyle  = 2 $ y de que $\displaystyle \Gamma = \Gamma_{1} \cup\Gamma_{2} $ sea cerrada tenemos que $\displaystyle \int _{\Gamma }f  $ no depende del punto inicial o final. En efecto,
\[\int _{\Gamma }f =\int _{\Gamma_{1}}f +\int _{\Gamma_{2}}f =  \int _{\Gamma_{2}}f +\int _{\Gamma_{1}}f  .\]
Entonces, si $\displaystyle \gamma, \omega : \mathbb{S}^{1} \to \C $ son equivalentes 
\[\int _{\gamma }f =\int _{\sigma } f .\]
En efecto, si $\displaystyle z = \gamma\left(1\right) $ y $\displaystyle \omega = \sigma\left(1\right) $ son distintos, 
\[\int _{\gamma }f =\int _{\Gamma_{1}}f +\int _{\Gamma_{2}}f =\int _{\Gamma_{2}}f +\int _{\Gamma_{1}}f =\int _{\sigma }f ,\]
donde $\displaystyle \Gamma_{1} $ denota el subarco de $\displaystyle \gamma\left(\mathbb{S}^{1}\right)=\sigma\left(\mathbb{S}^{1}\right) $ que comienza en $\displaystyle z $ y acaba en $\displaystyle w $ y se recorre en el mismo sentido que lo hacen $\displaystyle \gamma  $ y $\displaystyle \sigma  $; y $\displaystyle \Gamma_{2} $ es el subarco que comienza en $\displaystyle w $ y acaba en $\displaystyle z $ y se recorre en el mismo sentido que $\displaystyle \gamma  $ y $\displaystyle \sigma  $. 
\begin{definition}[Función primitiva]
Dado $\displaystyle \Omega \subset \C $ abierto y $\displaystyle f : \Omega \to \C $, decimos que $\displaystyle F : \Omega \to \C $ con $\displaystyle F \in \mathcal{H}\left(\Omega \right) $ es una \textbf{primitiva} de $\displaystyle f$ si $\displaystyle F'\left(z\right)=f\left(z\right) $, $\displaystyle \forall z \in \Omega  $. 
\end{definition}
\begin{observation}
Si $\displaystyle \Omega \subset \C $ es abierto conexo y $\displaystyle F_{1} $ y $\displaystyle F_{2} $ son primitivas de $\displaystyle f $, debe ser que $\displaystyle F_{1} $ y $\displaystyle F_{2} $ son iguales salvo quizás una constante.  
\end{observation}
\begin{theorem}
Sea $\displaystyle \Omega\subset \C $ abierto, $\displaystyle f \in \mathcal{C}\left(\Omega \right) $ y $\displaystyle F $ una primitiva de $\displaystyle f $ en $\displaystyle \Omega  $. Sea $\displaystyle \Gamma $ una curva de clase $\displaystyle \mathcal{C}^{1} $ a trozos que empieza en $\displaystyle \omega_{1} $ y termina en $\displaystyle \omega_{2} $. Entonces,
\[\int _{\Gamma }f =F\left(\omega_{2}\right)-F\left(\omega_{1}\right) .\]
\end{theorem}
\begin{proof}
	Supongamos que $\displaystyle \Gamma $ es $\displaystyle \mathcal{C}^{1} $. Sea $\displaystyle \gamma : \left[a,b\right] \to \Omega $ una parametrización de $\displaystyle \Gamma  $. Ahora, aplicando la regla de la cadena compleja 
	\[
	\begin{split}
	\int _{\gamma }f =\int^{b}_{a} f\left(\gamma\left(t\right)\right)\gamma'\left(t\right) \; dt= \int^{b}_{a} \frac{d}{dt}F\left(\gamma\left(t\right)\right) \; dt = F\left(\gamma\left(b\right)\right)-F\left(\gamma\left(a\right)\right)=F\left(\omega_{2}\right)-F\left(\omega_{1}\right) .
	\end{split}
	\]
En el caso de que $\displaystyle \Gamma $ sea $\displaystyle \mathcal{C}^{1} $ a trozos tendremos que 
\[\int _{\Gamma}f =\sum^{n}_{j = 1}\int _{\Gamma }f =\sum^{n}_{j = 1}\left(F\left(\omega_{2}^{j}\right)-F\left(\omega^{j}_{1}\right)\right) = F\left(\omega_{2}^{n}\right)-F\left(\omega_{1}^{1}\right) .\]
\end{proof}
\begin{colorary}
Si $\displaystyle \Gamma $ es una curva cerrada $\displaystyle \mathcal{C}^{1} $ a trozos y $\displaystyle f : \Omega\subset \C \to \C $, $\displaystyle f \in \mathcal{C}\left(\Omega \right) $ que admite primitiva en $\displaystyle \Omega  $, entonces $\displaystyle \int _{\Gamma }f =0 $. 
\end{colorary}
\begin{eg}
	La función $\displaystyle f\left(z\right)=\frac{1}{z} $ no admite primitiva en $\displaystyle \C\backslash \left\{ 0\right\}  $. En efecto, por el corolario anterior tenemos que 
	\[\int _{ \left\{ \left|z\right| = 1\right\} }\frac{1}{z} \; dz = \int^{2\pi }_{0} \frac{1}{e^{it}}ie^{it} \; dt =2\pi i .\]
\end{eg}
\begin{theorem}[Derivación bajo el signo de la integral]
Sean $\displaystyle \Omega\subset \C $ abierto, $\displaystyle \Gamma $ una curva orientada de clase $\displaystyle \mathcal{C}^{1} $ a trozos (no pedimos que esté en $\displaystyle \Omega  $). Sea $\displaystyle \varphi : \Gamma \times \Omega \to \C $ continua tal que $\displaystyle \forall \xi \in \Gamma  $ la aplicación $\displaystyle z \to \varphi\left(\xi, z\right)\in \mathcal{H}\left(\Omega \right) $ y tal que $\displaystyle \left(\xi, z\right)\to \frac{\partial \varphi}{\partial z}\left(\xi, z\right) $ es continua en $\displaystyle \Gamma \times \Omega  $.
Entonces
\[g\left(z\right):= \int _{\Gamma}\varphi\left(\xi, z\right) \; d\xi ,\]
es holomorfa en $\displaystyle z \in \Omega  $ y 
\[g'\left(z\right)=\int _{\Gamma}\frac{\partial\varphi}{\partial z}\left(\xi, z\right) \; d\xi .\]
\end{theorem}
\begin{proof}
	Sea $\displaystyle z_{0} \in \Omega  $ y $\displaystyle r > 0 $ tal que $\displaystyle \overline{B}\left(z_{0}, r\right)\subset \Omega  $. Para $\displaystyle \left|h\right|<r $, considerando el segmento $\displaystyle S = \left[z_{0}, z_{0} + h\right] \subset \Omega  $ parametrizado por $\displaystyle S_{h}\left(t\right)=z_{0} + th $, $\displaystyle t \in \left[0,1\right]  $, y aplicando el teorema anterior tenemos que 
	\[\varphi\left(\xi, z_{0} + h\right)-\varphi\left(\xi, z_{0}\right)=\int _{S_{h}}\frac{\partial\varphi}{\partial z}\left(\xi, z\right) \; dz = \int^{1}_{0} \frac{\partial \varphi}{\partial z}\left(\xi, S_{h}\right)S'_{h}\left(t\right) \; dt = \int^{1}_{0} \frac{\partial\varphi}{\partial z}\left(\xi,z_{0}+th\right)h \; dt \]
	para cada $\displaystyle \xi \in \Gamma $. Suponiendo que $\displaystyle \gamma : \left[a,b\right] \to \Omega  $ es una parametrización de $\displaystyle \Gamma $ tenemos que
	\[
	\begin{split}
		&\lim_{h \to 0}\frac{g\left(z_{0} + h\right)-g\left(z_{0}\right)}{h} =  \lim_{h \to 0}\frac{1}{h}\int _{\Gamma}\left(\varphi\left(\xi, z_{0} + h\right)-\varphi\left(\xi, z_{0}\right)\right) \; d\xi \\
		= &  \lim_{h \to 0}\frac{1}{h}\int _{\Gamma }\left(\int^{1}_{0} \frac{\partial\varphi}{\partial z}\left(\xi, z_{0} + sh\right)h \; d s\right) \; d\xi 
		=  \lim_{h \to 0}\int^{b}_{a} \left(\int^{1}_{0} \frac{\partial \varphi}{\partial z}\left(\gamma\left(t\right), z_{0} + sh\right) \; d s\right)\gamma'\left(t\right) \; dt \\
		= &  \lim_{h \to 0}\int _{\left[0,1\right] \times\left[a,b\right] }\frac{\partial \varphi}{\partial z}\left(\gamma\left(t\right), z_{0} + sh\right)\gamma'\left(t\right) \; d s \; dt  = \int _{\left[0,1\right] \times\left[a,b\right] }\lim_{h \to 0}\frac{\partial \varphi}{\partial z}\left(\gamma \left(t\right), z_{0} + sh\right)\gamma'\left(t\right) \; d s dt\\
		= & \int _{\left[0,1\right] \times\left[a,b\right] }\frac{\partial \varphi}{\partial z}\left(\gamma\left(t\right), z_{0}\right)\gamma'\left(t\right) \; d s \; dt = \int^{b}_{a} \frac{\partial \varphi}{\partial z}\left(\gamma\left(t\right), z_{0}\right)\gamma'\left(t\right) \; dt = \int _{\Gamma}\frac{\partial \varphi}{\partial z}\left(\xi, z\right) \; d\xi.
	\end{split}
	\]
En la cuarta igualdad hemos usado el teorema de Fubini y en la quinta el intercambio de la integral y el límite se justifica porque al ser $\displaystyle \frac{\partial \varphi}{\partial z} $ continua en el compacto $\displaystyle \Gamma \times \overline{B}\left(z_{0}, r\right) $ es uniformemente continua en este conjunto y por tanto
\[\lim_{h \to 0}\frac{\partial \varphi}{\partial z}\left(\gamma\left(t\right), z_{0} + sh\right)\gamma'\left(t\right)=\frac{\partial \varphi}{\partial z}\left(\gamma\left(t\right), z_{0}\right)\gamma'\left(t\right)\]
uniformemente en $\displaystyle \left(s,t\right)\in \left[0,1\right] \times\left[a,b\right]  $. 
\end{proof}
\begin{colorary}
Sean $\displaystyle z_{0} \in \C $ y $\displaystyle r > 0 $. Para todo $\displaystyle z \in D\left(z_{0}, r\right) $ es
\[\int _{ \left\{ \left|z-z_{0}\right|=r\right\} }\frac{1}{\xi-z} \; d\xi = 2\pi i ,\]
con orientación positiva. 
\end{colorary}
\begin{proof}
Tomamos 
\[\varphi: \partial D\left(z_{0}, r\right)\times \C \backslash \partial D\left(z_{0}, r\right)\to \C : \left(\xi, z \right)\to \frac{1}{\xi - z} .\]
Por construcción el denominador nunca se anula. En el dominio, la función $\displaystyle \varphi $ es continua y su derivada parcial es $\displaystyle \frac{\partial \varphi}{\partial z}=\frac{1}{\left(\xi-z\right)^{2}} $ que también es continua en el dominio. 
Así, la función
\[f\left(z\right):= \int _{ \left\{ \left|z-z_{0}\right|=r\right\} }\frac{1}{\xi-z} \; d\xi ,\]
está bien definida y por el teorema anterior es holomorfa en $\displaystyle D\left(z_{0}, r\right) $ con derivada
\[f'\left(z\right)=\int _{ \left\{ \left|z-z_{0}\right|=r\right\} }\frac{\partial \varphi}{\partial z} \; d\xi = \int _{ \left\{ \left|z-z_{0}\right|=r\right\} }\frac{1}{ \left(\xi-z\right)^{2}} \; d\xi .\]
Esta integral es nula para cada $\displaystyle z \in D\left(z_{0}, r\right) $ puesto que es la integral en una curva cerrada de una función que es continua y tiene primitiva $\displaystyle -\frac{1}{\xi-z} $. Se sigue que $\displaystyle f $ es constante en $\displaystyle D\left(z_{0}, r\right) $, por ser este un abierto conexo. Por ende, $\displaystyle \forall z \in D\left(z_{0}, r\right) $,
\[f\left(z\right)=f\left(z_{0}\right)=\int _{ \left\{ \left|z-z_{0}\right|=r\right\} }\frac{1}{\xi-z_{0}} \; d\xi =\int^{2\pi }_{0} \frac{rie^{it}}{re^{it}} \; dt = 2\pi i .\]
Hemos usdo la parametrización $\displaystyle \gamma\left(t\right)=z_{0} + re^{it} $ con $\displaystyle t \in \left[0,2\pi \right]  $. 
\end{proof}
\section{Teorema de Cauchy-Goursat}
\begin{theorem}[Teorema de Goursat]
	Sean $\displaystyle \Omega \subset \C $ un abierto, $\displaystyle T \subset \Omega  $ el borde de un triángulo con su interior contenido en $\displaystyle \Omega  $. Sea $\displaystyle f \in \mathcal{H}\left(\Omega \backslash \left\{ p\right\} \right) $ con $\displaystyle p \in \Omega  $. Entonces, $\displaystyle \int _{T}f =0 $.
\end{theorem}
\begin{proof}
	Sea $\displaystyle T = T_{0} $ el triángulo dado, con vértices $\displaystyle a,b,c \in \C $ y sea $\displaystyle \mathcal{T} = \mathcal{T}_{0} $ la unión de $\displaystyle T $ con la región interior a $\displaystyle T $ con borden $\displaystyle \partial \mathcal{T}=T $. Podemos suponer que $\displaystyle T $ está orientado positivamente y que sus lados se recorren en este orden: $\displaystyle \left[a,b\right]  $, $\displaystyle \left[b,c\right]  $, $\displaystyle \left[c,a\right]  $. Podemos distinguir varios casos:
	\begin{description}
		\item[Caso 1.] Supongamos que $\displaystyle p\not\in \mathcal{T} $. Uniendo los puntos medios de los lados $\displaystyle \left[a,b\right]  $, $\displaystyle \left[b,c\right]  $ y $\displaystyle \left[c,a\right]  $ obtenemos un triángulo que denotamos $\displaystyle T^{4}_{0} $ que divide $\displaystyle \mathcal{T}^{0}_{0} $ en cuatro subtriángulos sólidos $\displaystyle \mathcal{T}^{j}_{0} $ con bordes $\displaystyle T^{j}_{0} $ con $\displaystyle j = 1, \ldots, 4 $, cuyos perímetros son la mitad del perímetro del triángulo original $\displaystyle T $. 
% añadir imagen
Elijamos en los triángulos $\displaystyle T^{j}_{0} $ con $\displaystyle 1 \leq j \leq 3 $ las orientaciones inducidas por la orientación de $\displaystyle T $, y orientemos cada lado de $\displaystyle T^{4}_{0} $ de forma opuesta a la del mismo lado considerado como lado de $\displaystyle T^{j}_{0} $ para algún $\displaystyle j = 1, 2, 3 $ (equivalentemente, orientamos todos los triángulos $\displaystyle T^{j}_{0} $ de forma positiva). Entonces, tenemos que 
\[\int _{T}f =\sum^{4}_{j=1}\int _{T^{j}_{0}}f  ,\]
puesto que las integrales de la suma de la derecha se cancelan mutuamente en cada lado de $\displaystyle T^{4}_{0} $. Debe existir algún $\displaystyle j \in \left\{ 1, \ldots, 4\right\}  $ tal que 
\[ \left|\int _{T_{0}^{j}}f \right|\geq \frac{1}{4} \left|\int _{T_{0}}f \right| .\]
En caso contrario tendríamos que 
\[ \left|\int _{T_{0}}f \right|= \left|\sum^{4}_{j=1}\int _{T_{0}^{j}}f \right|<\sum^{4}_{j=1}\frac{1}{4} \left|\int _{T_{0}}f \right|= \left|\int _{T_{0}}f \right| .\]
Al $\displaystyle T^{j}_{0} $ que hemos elegido lo llamamos $\displaystyle T_{1} $, denotando por $\displaystyle \mathcal{T}_{1} $ la región interior a $\displaystyle T_{1} $ junto con el propio $\displaystyle T_{1} $. Repetimos el argumento anterior cambiando $\displaystyle T_{0} $ por $\displaystyle T_{1} $. Dividimos $\displaystyle \mathcal{T}_{1} $ en cuatro subtriángulos sólidos $\displaystyle \mathcal{T}^{j}_{1} $ con bordes $\displaystyle T^{j}_{1} $ con $\displaystyle j = 1, \ldots, 4 $, cuyos perímetros serán un cuerto del perímetro del triángulo original $\displaystyle T $. 
Elegimos uno de los cuatro triángulos $\displaystyle T^{j}_{1} $ para el que se tenga 
\[ \left|\int _{T^{j}_{1}}f \right|\geq \frac{1}{4} \left|\int _{T_{1}}f \right| .\]
De esta forma, 
\[ \left|\int _{T^{j}_{1}}f \right|\geq \frac{1}{4^{2}} \left|\int _{T}f \right| .\]
Reiterando este proceso obtenemos por inducción una sucesión $\displaystyle \left\{ \mathcal{T}_{j}\right\} _{j \in \N} $ de triángulos sólidos con bordes $\displaystyle T_{j} $ tales que 
\begin{enumerate}
\item $\displaystyle \mathcal{T}_{j}\subset\mathcal{T}_{j-1} $, $\displaystyle \forall j \in \N $ con $\displaystyle \mathcal{T}_{0}=\mathcal{T} $. 
\item $\displaystyle \left|\int _{T}f \right|\leq 4^{j} \left|\int _{T_{j}}f \right| $. 
\item $\displaystyle 2^{j}\ell\left(T_{j}\right)=\ell\left(T\right) $. 
\end{enumerate}
En particular, como la longitud de un triángulo es mayor o igual que su diámetro, se tiene que $\displaystyle \left\{ \mathcal{T}_{j}\right\} _{j \in \N} $ es una familia de compactos en $\displaystyle \C \equiv \R^{2} $ cuyos diámetros tienden a cero, por lo que existe un único $\displaystyle z_{0} \in \C $ tal que 
\[z_{0}\in \bigcap_{j \in \N}\mathcal{T}_{j} .\]
Puesto que estamos suponiendo que $\displaystyle p \not\in \mathcal{T} $, la función $\displaystyle f $ es holomorfa en $\displaystyle z_{0} $ y por tanto, dado cualquier $\displaystyle \epsilon > 0 $ existe $\displaystyle \delta > 0 $ tal que si $\displaystyle \left|z-z_{0}\right|\leq \delta  $, entonces
\[ \left|f\left(z\right)-f\left(z_{0}\right)-f'\left(z_{0}\right)\left(z-z_{0}\right)\right|\leq \epsilon \left|z-z_{0}\right| .\]
Por otro lado, como $\displaystyle \lim_{j \to \infty}\ell\left(T_{j}\right)=0 $, existe $\displaystyle j_{0} \in \N $ tal que si $\displaystyle j \geq j_{0} $ entonces
\[diam\left(T_{j}\right)\leq \ell\left(T_{j}\right)<\delta ,\]
luego $\displaystyle \left|z-z_{0}\right|\leq \delta  $ para todo $\displaystyle z \in T_{j} $, de donde deducimos que para $\displaystyle j \geq j_{0} $, 
\[
\begin{split}
	\left|\int _{T_{j}}f \right|= & \left|\int _{T_{j}}f\left(z\right)-f\left(z_{0}\right)-f'\left(z_{0}\right)\left(z-z_{0}\right) \; dz\right|\\
	\leq& \epsilon \sup_{z\in T_{j}} \left|z-z_{0}\right| \ell\left(T_{j}\right) \leq \epsilon \ell\left(T_{j}\right)^{2}=\frac{\epsilon }{4^{j}}\ell\left(T\right)^{2},
\end{split}
\]
donde en la primera igualdad hemos usado que 
\[\int _{T_{j}}f\left(z_{0}\right)+f'\left(z_{0}\right)\left(z-z_{0}\right) \; dz = 0 ,\]
por ser la integral en una curva cerrada de una función que tiene primitiva. En la segunda desigualdad hemos vuelto a usar que el diámetro de un triángulo es menor o igual que su perímetro, y en la última igualdad hemos utilizado la tercera propiedad de la sucesión de triángulos. Combinando la desigualdad anterior con la propiedad \textbf{(2)} de la sucesión de triángulos se obtiene que 
\[ \frac{1}{4^{j}} \left|\int _{T}f \right|\leq \left|\int _{T_{j}}f \right|\leq \frac{\epsilon }{4^{j}}\ell\left(T\right)^{2} \Rightarrow \left|\int _{T}f \right|\leq \epsilon \ell\left(T\right)^{2} .\]
Como es cierto para todo $\displaystyle \epsilon > 0 $ debe ser que $\displaystyle \int _{T}f =0 $. 
\item[Caso 2.] Supongamos que $\displaystyle p \in \mathcal{T} $. Podemos distinguir tres casos. 
	\begin{description}
	\item[Caso 2.1.] Si $\displaystyle p $ es un vértice de $\displaystyle T $ podemos suponer sin pérdida de generalidad que $\displaystyle p = a $. Como $\displaystyle f $ es continua en $\displaystyle T $ compacto, existe $\displaystyle M > 0 $ tal que $\displaystyle \left|f\left(z\right)\right|\leq M $, $\displaystyle \forall z \in T $. 
		Dado $\displaystyle \epsilon > 0 $, consideremos los puntos $\displaystyle x \in \left[a,b\right]  $ e $\displaystyle y \in \left[c,a\right]  $ tales que la longitud del triángulo $\displaystyle T_{1} $ de vértices $\displaystyle a,x,y $ sea menor que $\displaystyle \frac{\epsilon }{M} $. 
		Denotamos por $\displaystyle T_{2} $ el triángulo de vértices $\displaystyle x,b,y $ y por $\displaystyle \mathcal{T}_{3} $ el de vértices $\displaystyle b,c,y $ con las orientaciones inducidas por la de $\displaystyle T $. Por lo demostrado en el \textbf{caso 1} tenemos que 
		\[\int _{T_{2}}f =\int _{T_{3}}f =0 .\]
		Por tanto, 
		\[ \left|\int _{T}f \right|= \left|\sum^{3}_{j=1}\int _{T_{j}}f \right|= \left|\int _{T_{1}}f \right|\leq M \ell\left(T_{1}\right) \leq \epsilon  .\]
		Como esto es cierto para todo $\displaystyle \epsilon > 0 $ debe ser que $\displaystyle \int _{T}f =0 $. 
	\item[Caso 2.2.] Supongamos que $\displaystyle p $ está en un lado de $\displaystyle T $. Digamos por ejemplo que $\displaystyle p \in \left[a,b\right]  $. Sea $\displaystyle T_{1} $ el triángulo de vértices $\displaystyle a,p,c $ y $\displaystyle T_{2} $ el triángulo de vértices $\displaystyle p,b,c $. Usando el caso anterior tenemos que 
		\[\int _{T}f =\int _{T_{1}}f +\int _{T_{2}}f = 0 + 0 = 0 .\]
	\item[Caso 2.3.] Supongamos que $\displaystyle p $ está en el interior de $\displaystyle \mathcal{T} $. Definimos $\displaystyle T_{1} $ como el triángulo de vértices $\displaystyle a,p,c $; $\displaystyle T_{2} $ el de vértices $\displaystyle a,b,p $; y $\displaystyle T_{3} $ el de $\displaystyle b,c,p $. 
		Usando el \textbf{caso 2.1} tenemos que 
		\[\int _{T}f =\int _{T_{1}}f +\int _{T_{2}}f +\int _{T_{3}}f =0 + 0 + 0 = 0 .\]
	\end{description}
	\end{description}
\end{proof}

\begin{proof}
Procedemos por inducción. Sea $\displaystyle T_{0} $ el borde del triángulo inicial con vértices $\displaystyle a,b,c \in \C $. Sea $\displaystyle S_{0} $ el triángulo sólido cuyo borde es $\displaystyle \partial S_{0} = T_{0} $. Distinguimos varios casos:
\begin{itemize}
	\item Supongamos que $\displaystyle p \not\in S_{0} $. Uniendo los puntos medios de los lados $\displaystyle \left[a,b\right]  $, $\displaystyle \left[b,c\right]  $ y $\displaystyle \left[c,a\right]  $ obtenemos un triángulo que denotamos $\displaystyle T_{0}^{4} $ que divide a su vez $\displaystyle T_{0} $ en cuatro subtriángulos sólidos $\displaystyle T^{j}_{0} $, cuyos perímetros son la mitad del perímetro del triángulo original $\displaystyle T_{0} $. Elegimos en los triángulos $\displaystyle T^{j}_{0} $ con $\displaystyle 1 \leq j \leq 3 $, las orientaciones inducidas por la orientación de $\displaystyle T_{0} $. Entonces, tenemos que
		\[\int _{T}f =\int _{T_{0}}f =\sum^{4}_{j = 1}\int _{T^{j}_{0}}f .\]
		Por el principio del palomar debe existir $\displaystyle j \in \left\{ 1, \ldots, 4\right\}  $ con 
		\[ \left|\int _{T^{j}_{0}}f \right|\geq \frac{1}{4} \left|\int _{T}f \right| .\]
	Si no fuese así tendríamos que 
	\[ \left|\int _{T_{0}}f \right|= \left|\sum^{4}_{j = 1}\int _{T^{j}_{0}}f \right|<\sum^{4}_{j = 1}\frac{1}{4} \left|\int _{T_{0}}f \right| = \left|\int _{T_{0}}f \right| ,\]
	que es una contradicción. Si cogemos ahora $\displaystyle T_{1} = T_{0}^{j} $ con el $\displaystyle j $ que hemos cogido antes y sea $\displaystyle S_{1}  $ el triángulo sólido tal que $\displaystyle \partial S_{1}= T_{1} $. Procediendo inductivamente, se obtiene una sucesión de triángulos sólidos $\displaystyle \left\{ S_{j}\right\} _{j \geq 0} $ con bordes $\displaystyle \partial S_{j}=\partial T_{j} $ que satisface las siguientes propiedades
	\begin{itemize}
	\item $\displaystyle \left|\int _{T_{j}}f \right| \geq \frac{1}{4^{j}} \left|\int _{T}f \right| $, $\displaystyle \forall j \geq 0 $. 
	\item $\displaystyle S_{j + 1}\subset S_{j} $, $\displaystyle \forall j\geq 0 $. 
	\item $\displaystyle \ell\left(T_{j}\right)=\frac{1}{2^{j}}\ell\left(T\right) $, $\displaystyle \forall j \geq 0 $. 
	\end{itemize}
	Además, $\displaystyle diam\left(T_{j}\right)\to 0 $ si $\displaystyle j \to \infty $. Por tanto, 
	\[\bigcap_{j \geq 0}S_{j}= \left\{ z_{0}\right\}  ,\]
	para algún $\displaystyle z_{0}\in \overline{B}\left(x_{0}, r\right) $ con $\displaystyle \overline{B}\left(x_{0}, r\right)\subset \Omega  $. Por construcción, 
	\[diam\left(\bigcap_{j \geq 0}S_{j}\right)\leq diam S_{j} \xrightarrow{j \to \infty}0 .\]
	Así, $\displaystyle diam \left(\bigcap_{j \geq 0}S_{j}\right) $ tiene a lo sumo un punto. Como $\displaystyle \bigcap_{ j \geq 0}S_{j} $ es una intersección de compactos encajados, tendremos que no es el vacío. Por ser $\displaystyle f \in \mathcal{H}\left(\Omega \backslash \left\{ p\right\} \right) $ entonces $\displaystyle f $ es holomorfa en $\displaystyle z_{0} \in S $ por lo que $\displaystyle z _{0}\neq p $. Para $\displaystyle \epsilon > 0 $, existe $\displaystyle \delta > 0 $ tal que si $\displaystyle \left|z-z_{0}\right|<\delta  $ entonces
	\[ \left|f\left(z\right)-f\left(z_{0}\right)-f'\left(z_{0}\right)\left(z-z_{0}\right)\right|\leq \epsilon \left|z-z_{0}\right| .\]
	Ahora, existe $\displaystyle j_{0} \in \N $ tal que $\displaystyle \ell\left(T_{j}\right)<\delta  $, entonces para $\displaystyle \left|z-z_{0}\right|<\delta  $, $\displaystyle \forall z \in T_{j} $, $\displaystyle \forall j \geq j_{0} $, entonces
	\[ \left|\int _{T_{j}}f \right|= \left|\int _{T_{j}}f\left(z\right)-f\left(z_{0}\right)-f'\left(z_{0}\right)\left(z-z_{0}\right) \right|\leq \epsilon \sup_{z\in T_{j}} \left|z-z_{0}\right|\ell\left(T_{j}\right) \leq \epsilon \ell\left(T_{j}\right)\frac{\ell\left(T\right)}{2^{j}}\leq\frac{\epsilon }{4^{j}}\ell\left(T\right) .\]
Por otro lado, 
\[ \left|\int _{T}f \right| \leq 4^{j} \left|\int _{T_{j}}f \right| \leq \epsilon \ell\left(T\right) .\]
Así, tendremos que $\displaystyle \int _{T}f =0 $. 
\item Supongamos que $\displaystyle p \in S $. Distinguimos varios casos.
	\begin{itemize}
	\item Si $\displaystyle p $ es un vértice, podemos suponer sin pérdida de generalidad que $\displaystyle p = a $. 
	\end{itemize}
	
\end{itemize}
\end{proof}

